\section{Implementación}
\label{sec:implementacion}

\subsection{Stack tecnológico}

\begin{table}[H]
\centering
\small
\begin{tabular}{p{3cm}p{9cm}}
\toprule
\textbf{Capa} & \textbf{Tecnología y versión} \\
\midrule
Frontend & Angular 21.2.12, servido en producción por nginx (build multi-stage, ver \texttt{Frontend/Dockerfile}) \\
Backend & Spring Boot 3.2.1 sobre Java 17 (compilado con JDK 21 en modo \texttt{--release 17}) \\
Persistencia & PostgreSQL 15 (Flyway para migraciones versionadas) \\
Caché & Redis 7, vía \texttt{spring-boot-starter-data-redis} \\
Autenticación & JJWT 0.12.5, JWT de 7 claims RFC~7519, refresh token con blacklist en Redis \\
Documentación API & Springdoc OpenAPI 2.3.0 (Swagger UI), versionado dinámico \texttt{/api/v1/} \\
Generación de PDF & iText (actas y reportes) \\
CI/CD & GitHub Actions (\texttt{.github/workflows/ci.yml}, \texttt{backup.yml}) \\
\bottomrule
\end{tabular}
\caption{The system's real technology stack.}
\end{table}

El workflow \texttt{.github/workflows/ci.yml} corre en cada \textit{push}, levantando Postgres y Redis
reales (no simulados) para ejecutar build, la suite de tests con JaCoCo y el análisis estático de
seguridad del backend. El Anexo C (Sección~\ref{anexo:ci}) documenta tres corridas verdes reales
verificadas contra la API pública de GitHub, con sus datos crudos versionados en el repositorio.

\subsection{Seguridad implementada}

La autenticación usa JWT con los 7 claims de RFC~7519 (\texttt{jti}, \texttt{iss}, \texttt{sub},
\texttt{aud}, \texttt{iat}, \texttt{nbf}, \texttt{exp}), refresh token con rotación y blacklist
almacenados en Redis, rate limiting en \texttt{/api/v1/auth/login} (6 intentos/60s $\rightarrow$
HTTP~429), y cabeceras de seguridad HTTP (HSTS, CSP, X-Frame-Options, X-Content-Type-Options)
configuradas tanto en Spring Security como en nginx. El control de acceso usa
\texttt{@PreAuthorize} con un sistema de permisos dinámico (\texttt{presus.permisos}/\texttt{rol\_permisos},
editable sin recompilar) además de los cuatro roles fijos (ADMIN, DOCENTE, COORDINADOR, ESTUDIANTE). Los
procedimientos almacenados y las consultas JPA son parametrizados; el análisis estático con find-sec-bugs
confirmó 0 hallazgos de la regla \texttt{SQL\_NONCONSTANT\_STRING\_PASSED\_TO\_EXECUTE}
(Sección~\ref{sec:evaluacion}). Más allá del análisis automatizado, una revisión manual dirigida del
control de acceso \emph{por recurso} (¿puede este usuario ver/modificar \emph{este} registro en
particular, no solo llamar a este endpoint?) encontró y corrigió 6 fallos reales de autorización --- IDOR
de lectura y de escritura, un endpoint administrativo sin el permiso correspondiente, y mass-assignment
en la creación de usuarios --- detallados con su verificación en vivo en
Sección~\ref{sec:seguridad-corregida}.

\subsection{Integración de API externa y caché}

El servicio \texttt{ExternalApiServiceImpl} consume la API pública de universidades de Hipo Labs vía
\texttt{WebClient} reactivo, con reintentos con backoff exponencial (3 intentos) ante fallos de red, y
resultado cacheado en Redis (\texttt{@Cacheable}) con TTL de 10 minutos para universidades y 5 minutos
para solicitudes. Este es el componente cuyo efecto de caché se mide cuantitativamente en la
Sección~\ref{sec:evaluacion} (RQ2).

\subsection{Despliegue}

El sistema se ejecuta localmente con \texttt{docker compose up -d --build} (cuatro contenedores:
postgres, redis, backend, nginx+frontend). Para producción, se preparó --- pero a la fecha de este
informe \textbf{no se ha completado} --- un despliegue en Railway (\texttt{docs/despliegue/}), con un
servicio Railway por contenedor, comunicados por red privada, y HTTPS gestionado automáticamente por el
proveedor. La preparación incluyó: un \texttt{Frontend/Dockerfile} nuevo (antes no existía, el frontend
solo se montaba como volumen local), la corrección de un \texttt{package-lock.json} desincronizado que
rompía \texttt{npm ci} en un entorno Linux limpio, y la habilitación del detalle completo de
\texttt{/actuator/health} (antes solo devolvía \texttt{\{"status":"UP"\}} sin desglose de componentes).

\subsection{Usuario de demostración}

Se sembró un usuario de demostración (\texttt{demo@uteq.edu.ec}, rol Coordinador) además del
administrador, para que un evaluador externo pueda explorar el flujo académico completo sin necesidad de
registrarse. Las credenciales de ambos usuarios se entregan al tribunal junto con el enlace de despliegue
en el momento de la evaluación; se retiraron explícitamente de \texttt{README.md} durante este mismo
proceso de correcciones para no dejarlas expuestas en texto plano en el historial público del
repositorio.

\subsection{Capturas reales del sistema}

Las Figuras~\ref{fig:cap-solicitud}, \ref{fig:cap-anteproyecto}, \ref{fig:cap-tramites} y
\ref{fig:cap-horario} muestran pantallas reales del sistema en producción (Angular 21, sobre nginx,
build multi-stage), en el orden del flujo real de un estudiante: crear la solicitud, cargar el
anteproyecto, dar seguimiento al trámite y consultar el horario asignado.
Son los mismos archivos PNG versionados en \texttt{Frontend/public/img/} desde el commit inicial de la
estructura del proyecto.

\begin{figure}[H]
\centering
\includegraphics[width=0.70\textwidth]{figuras/cap-nueva-solicitud}
\caption{Módulo ``Nueva Solicitud'' --- el estudiante registra su solicitud de pre-sustentación
(fuente: \texttt{Frontend/public/img/NuevaSolicitud.png}).}
\label{fig:cap-solicitud}
\end{figure}

\begin{figure}[H]
\centering
\includegraphics[width=0.70\textwidth]{figuras/cap-cargar-anteproyecto}
\caption{Módulo ``Cargar Anteproyecto'' --- carga del PDF del anteproyecto antes de la pre-sustentación
(fuente: \texttt{Frontend/public/img/CargarAnteproyecto.png}).}
\label{fig:cap-anteproyecto}
\end{figure}

\begin{figure}[H]
\centering
\includegraphics[width=0.70\textwidth]{figuras/cap-mis-tramites}
\caption{Módulo ``Mis Trámites'' --- vista de seguimiento del estado de la solicitud por parte del
estudiante (fuente: \texttt{Frontend/public/img/MisTramites.png}).}
\label{fig:cap-tramites}
\end{figure}

\begin{figure}[H]
\centering
\includegraphics[width=0.70\textwidth]{figuras/cap-mi-horario}
\caption{Módulo ``Mi Horario'' --- vista del cronograma de pre-sustentación asignado al estudiante
(fuente: \texttt{Frontend/public/img/MiHorario.png}).}
\label{fig:cap-horario}
\end{figure}

\subsection{Asistente virtual de ayuda (Chatbot)}

El sistema incluye un asistente de ayuda accesible desde cualquier pantalla autenticada
(\texttt{ChatbotController}, endpoint \texttt{POST /api/v1/chatbot/ask}). No es un modelo de lenguaje:
\texttt{ChatbotService} implementa un enrutador de intenciones por palabras clave sobre el mensaje del
usuario (\texttt{solicitud}, \texttt{anteproyecto}, \texttt{notificación}, \texttt{perfil},
\texttt{sustentación}, \texttt{contraseña}, \texttt{ayuda}), devolviendo una respuesta guía fija junto
con una lista de opciones sugeridas para cada tema; si ningún patrón coincide, responde con un mensaje
por defecto que remite a las preguntas frecuentes. La sesión se valida server-side
(\texttt{SecurityContextHolder}) antes de procesar cualquier mensaje: un usuario no autenticado recibe
una respuesta que le pide iniciar sesión, en vez de que la petición llegue a fallar sin explicación.

Fue, hasta la auditoría del 2026-09-05, el único controlador del proyecto sin ninguna anotación de
autorización (protegido solo por la regla global de \texttt{SecurityConfig}); se corrigió agregando
\texttt{@PreAuthorize("isAuthenticated()")} a nivel de clase, el mismo nivel que otros controladores de
consulta general cuya respuesta no depende del rol (Sección~\ref{sec:seguridad-corregida}). Tanto el
controlador (\texttt{ChatbotControllerTest}) como el servicio (\texttt{ChatbotServiceTest}) tienen
prueba automatizada dedicada, incluido el caso sin autenticar y la cobertura de cada rama de intención.

\textbf{Cierre de brecha (examen suspenso, P7):} las dos pruebas de arriba invocan el método Java del
controlador/servicio directamente con mocks, sin pasar por el \emph{dispatcher} real de Spring MVC ni
por la cadena de filtros de seguridad --- no ejercitaban el endpoint HTTP en sí. Se agregó
\texttt{ChatbotControllerIntegrationTest} (\texttt{@WebMvcTest} + \texttt{SecurityConfig} real, mismo
patrón que \texttt{MeControllerTest}), que hace peticiones HTTP reales contra
\texttt{POST /api/v1/chatbot/ask}: sin token (401), con token válido y cuerpo JSON serializado de
verdad (200, respuesta del mock del servicio visible en el JSON de salida), y con un cuerpo JSON mal
formado (400). Corre como parte de \texttt{./mvnw test} en el job \texttt{backend} de
\texttt{.github/workflows/ci.yml} igual que el resto de la suite, sin configuración adicional.
\newpage
